\section*{अध्याय \ref{ch:b3:product} --- गुणन माप,
फुबिनी, चर परिवर्तन}

\begin{solution}{exo:b3:product:1}
$\Delta$ $\intcc01^2$ में संवृत है, अतः बोरेल, और $\mathcal
B(\intcc01^2)$ गुणन $\sigma$-बीजगणित है (\cref{prop:b3:product:sections}(ख))।
एक ओर से पुनरावृत्ति करने पर:
\[
\int_{\intcc01}\Bigl(\int \mathbf 1_\Delta(x,y)
\,\dd\lambda(y)\Bigr)\dd\mu(x)
= \int \lambda(\{x\})\,\dd\mu(x) = 0 ;
\]
और दूसरी ओर से:
\[
\int_{\intcc01}\Bigl(\int\mathbf
1_\Delta(x,y)\,\dd\mu(x)\Bigr)\dd\lambda(y)
= \int \mu(\{y\})\,\dd\lambda(y) = \int 1\,\dd\lambda = 1 .
\]
जो परिकल्पना विफल होती है वह अगणनीय $\intcc01$ पर गणना \omterm{def:b3:measure:measure}{माप} $\mu$ की
$\sigma$-परिमितता है: परिमित-$\mu$ समुच्चयों का कोई गणनीय कुल उसे आच्छादित
नहीं करता।
\end{solution}

\begin{solution}{exo:b3:product:2}
$\intcc0A\times\intoo0{+\infty}$ पर: $\int_0^A\int_0^\infty\eu^{-xy}\abs{\sin x}\,\dd y\,\dd x =
\int_0^A\frac{\abs{\sin x}}x\dd x \leq A < \infty$ (निरपेक्ष मान के लिए टोनेली): अतः फुबिनी लागू होती
है, और चूँकि $\int_0^\infty\eu^{-xy}\dd y = \frac1x$,
\[
\int_0^A\frac{\sin x}x\dd x
= \int_0^\infty\Bigl(\int_0^A\eu^{-xy}\sin x\,\dd
x\Bigr)\dd y
= \int_0^\infty \frac{1 - \eu^{-Ay}(\cos A + y\sin A)}{1 +
y^2}\,\dd y
\]
(आंतरिक समाकल: $\operatorname{Im}\int_0^A\eu^{(\iu -
y)x}\dd x$, सीधे परिकलित)। $A \to \infty$ पर संशोधन पद $\int_0^\infty\eu^{-Ay}\frac{1 +
y}{1 + y^2}\dd y \leq \frac32\int_0^\infty\eu^{-Ay}\dd y =
\frac3{2A} \to 0$ से परिबद्ध
है; और मुख्य पद $\int_0^\infty\frac{\dd y}{1+y^2} = \frac\pi2$ है। अतः $\int_0^\infty\frac{\sin x}x\dd x = \frac\pi2$ --- अर्थात् फुबिनी से दिरिक्ले
का समाकल।
\end{solution}

\begin{solution}{exo:b3:product:3}
(क) परत-केक (\cref{prop:b3:product:layercake}): $\int
f\,\dd\mu = \int_0^\infty\mu(f > t)\,\dd t$, और $t \mapsto
\mu(f > t)$ अनवर्धमान है। $[n-1, n]$ पर: $\mu(f \geq n)
\leq \mu(f > t) \leq \mu(f > n - 1) \leq \mu(f \geq n - 1)$; इकाई
अंतरालों पर समाकलों का योग करने पर:
\[
\sum_{n\geq1}\mu(f \geq n) \leq \int f\,\dd\mu \leq
\sum_{n\geq1}\mu(f \geq n - 1) = \mu(f \geq 0) +
\sum_{n\geq1}\mu(f\geq n) \leq \mu(X) + \sum_{n\geq1}\mu(f\geq
n).
\]

(ख) (क) को $\abs f$ पर लगाइए: समाकल और श्रेणी की परिमितता तुल्य हैं
(अतिरिक्त $\mu(X)$ परिमित है)।
\end{solution}

\begin{solution}{exo:b3:product:4}
चूँकि $x \ne 0$ के लिए $f(x,y) = \partial_y\bigl(\frac{y}{x^2+y^2}\bigr)$:
\[
\int_0^1 f(x, y)\,\dd y = \frac{1}{x^2 + 1}
\ \Longrightarrow\
\int_0^1\!\!\int_0^1 f\,\dd y\,\dd x =
\int_0^1\frac{\dd x}{1 + x^2} = \frac\pi4 ;
\]
प्रतिसममिति $f(y,x) = -f(x,y)$ से दूसरा क्रम $-\frac\pi4$ देता है। निरपेक्ष मान: $0 < y < x$ के
लिए
\[
\int_0^x f(x,y)\,\dd y = \Bigl[\frac{y}{x^2 +
y^2}\Bigr]_0^x = \frac1{2x},
\quad\text{तथा वहाँ } f \geq 0 \text{, अतः}\quad
\int_0^1\!\!\int_0^1\abs f \geq \int_0^1\frac{\dd x}{2x} =
+\infty .
\]
फुबिनी के साथ कोई विरोधाभास नहीं: उसकी परिकल्पना $f \in L^1$ विफल हो जाती
है, और दोनों पुनरावृत्त समाकल केवल दो भिन्न संख्याएँ हैं।
\end{solution}

\begin{solution}{exo:b3:product:5}
(क) $(\mathbf 1_{\intcc01}*\mathbf 1_{\intcc01})(x) =
\lambda\bigl(\intcc01\cap\intcc{x-1}x\bigr)$: $x
\notin \intcc02$ के लिए $0$, $0 \leq x \leq 1$ के लिए $x$, $1
\leq x \leq 2$ के लिए $2 - x$:
अर्थात् तंबू। फिर से संवलन करने पर $\intcc03$ पर कोई $\mathcal
C^1$ खंडशः-द्विघातीय
उभार मिलता है (द्विघातीय बी-प्रत्यास्थक): प्रत्येक संवलन चिकनाई की
एक कोटि जोड़ देता है --- \cref{ch:b3:lp} के मृदुकारकों के पीछे यही मृदुकरण
सिद्धांत है।

(ख) यदि $x \notin \overline{\operatorname{supp}f +
\operatorname{supp}g}$ हो, तो $x$ के चारों ओर योग-समुच्चय से असंयुक्त कोई
गोला है; और $y \in \operatorname{supp}g$ के लिए $x - y
\notin\operatorname{supp}f$, अतः समाकल्य सर्वथा लुप्त हो जाता है:
$x$ के निकट $f * g = 0$।

(ग) $x_n \to x$ के लिए: $(f*g)(x_n) = \int
f(y)g(x_n - y)\,\dd y$; समाकल्य बिंदुशः अभिसरित होते हैं ($g$ की
\omterm{def:b3:topology:continuity}{सांतत्य}) और उन पर $\norm
g_\infty\,\abs f \in L^1$ का प्रभुत्व है: प्रभावी अभिसरण प्रमेय $(f*g)(x_n) \to (f*g)(x)$
दे देती है।
\end{solution}

\begin{solution}{exo:b3:product:6}
(क) फुबिनी और आगमन से, अंतिम निर्देशांक के अनुदिश काटते हुए:
\[
\lambda_d(\Delta_d) = \int_0^1
\lambda_{d-1}\bigl((1 - t)\,\Delta_{d-1}\bigr)\,\dd t
= \lambda_{d-1}(\Delta_{d-1})\int_0^1(1 - t)^{d-1}\dd t
= \frac{\lambda_{d-1}(\Delta_{d-1})}{d},
\]
जिसमें प्रसार नियम $\lambda_{d-1}(\rho A) =
\rho^{d-1}\lambda_{d-1}(A)$ (\cref{thm:b3:product:linearchange}) का उपयोग हुआ है; और $\lambda_1(\Delta_1) = 1$ के साथ:
आयतन $\frac1{d!}$।

(ख) $v_2 = \pi/\Gamma(2) = \pi$; $v_3 =
\pi^{3/2}/\Gamma(\frac52) = \pi^{3/2}/(\frac32\cdot\frac12
\sqrt\pi) = \frac{4\pi}3$। दीर्घवृत्तज $T = \operatorname{diag}(a_1, \dots, a_d)$ के साथ $T(B(0,1))$ है: \cref{thm:b3:product:linearchange} आयतन $v_d\prod a_i$ दे
देता है।
\end{solution}

\begin{solution}{exo:b3:product:7}
$\R^2$ में \omterm{ex:b3:product:polar}{ध्रुवीय निर्देशांक} (\cref{ex:b3:product:polar}):
\[
\int_{B(0,1)}\frac{\dd x\dd y}{(x^2+y^2)^s}
= 2\pi\int_0^1 r^{1 - 2s}\,\dd r,
\qquad
\int_{\R^2\setminus B(0,1)} = 2\pi\int_1^\infty r^{1-2s}\dd r:
\]
जो परिमित है तभी और केवल तभी जब $1 - 2s > -1$ ($s < 1$), क्रमशः $1 - 2s < -1$ ($s >
1$)।
$\R^d$ में गोलीय निर्देशांकों से बचकर परत-केक से: $\lambda_d(\{\norm x^{-2s} > t\}\cap B(0,1)) =
\lambda_d(B(0, \min(1, t^{-1/2s}))) = v_d\min(1, t^{-d/2s})$, और $\int_0^\infty v_d\min(1, t^{-d/(2s)})\dd t < \infty$ तभी
और केवल तभी जब $\frac d{2s} > 1$, अर्थात् $s < \frac d2$; और बाहरी समाकल अभिसरित होता है
तभी और केवल तभी जब $s > \frac d2$ (पूरक क्षेत्र पर वही परिकलन)।
\end{solution}

\begin{solution}{exo:b3:product:8}
टोनेली (धनात्मक समाकल्य) और चर परिवर्तन $(x, y) = \Phi(u, v) = (uv,\ u(1-v))$ से, जो $\intoo0\infty\times\intoo01$ का विवृत
चतुर्थांश पर कोई $\mathcal C^1$ अवकल समरूपता है और जिसमें
\[
\det D\Phi = \det\begin{pmatrix} v & u\\ 1 - v & -u
\end{pmatrix} = -uv - u(1 - v) = -u,
\qquad \abs{\det} = u :
\]
\[
\Gamma(p)\Gamma(q) =
\iint x^{p-1}y^{q-1}\eu^{-x-y}\dd x\,\dd y
= \iint (uv)^{p-1}\bigl(u(1{-}v)\bigr)^{q-1}\eu^{-u}\,u\,
\dd u\,\dd v
= \Gamma(p + q)\,B(p, q).
\]
$B(p,q)$ में $t = \sin^2\theta$ प्रतिस्थापित करने पर $2\int_0^{\pi/2}\sin^{2p-1}\theta\cos^{2q-1}\theta\,\dd\theta =
B(p, q)$ मिलता है। वालिस: $W_n =
\int_0^{\pi/2}\sin^n\theta\,\dd\theta = \frac12B\bigl(\frac{n +
1}2, \frac12\bigr) = \frac{\Gamma(\frac{n+1}2)\sqrt\pi}
{2\,\Gamma(\frac n2 + 1)}$ ---
उदाहरणार्थ $\Gamma(n + \frac12) = \frac{(2n)!}{4^nn!}\sqrt\pi$ का उपयोग करते हुए $W_{2n} =
\frac\pi2\cdot\frac{(2n)!}{4^n(n!)^2}$।
\end{solution}

\begin{solution}{exo:b3:product:9}
सूचक: $\int\mathbf 1_B\,\dd(T_*\mu) = T_*\mu(B) =
\mu(T^{-1}B) = \int\mathbf 1_B\circ T\,\dd\mu$; रैखिकता इसे सरल $g$ तक और एकदिष्ट अभिसरण प्रमेय $g \geq 0$
तक बढ़ा देती है --- यही अंतरण सूत्र है। यह विशुद्ध रूप से औपचारिक
है: वह $T_*\mu$ के सापेक्ष समाकलों को फिर से व्यक्त करता है पर इस बारे
में कुछ नहीं कहता कि $T_*\mu$ \emph{है क्या}। \cref{thm:b3:product:linearchange} और \cref{thm:b3:product:changeofvar} की अंतर्वस्तु
पहचान
\[
\Phi_*\bigl(\lambda_d\restriction_U\bigr) =
\abs{\det D\Phi^{-1}}\,\lambda_d\restriction_V
\quad\text{(एक घनत्व माप)},
\]
है, अर्थात् \omterm{def:b3:measure:lebesgueouter}{लेबेग माप} के अग्रसारण का परिकलन --- जिसका वैश्लेषिक
निवेश $\Phi$ की अवकल ज्यामिति है, न कि माप-सैद्धांतिक औपचारिकता।
\end{solution}

\begin{solution}{exo:b3:product:10}
टोनेली से गाउसीय गुणनखंडित हो जाता है, अतः $G_1 =
\int_\R\eu^{-s^2}\dd s = \sqrt\pi$ और $\int_\R
s^2\eu^{-s^2}\dd s = \frac{\sqrt\pi}2$ के साथ
(खंडशः समाकलन कीजिए):
\[
\int_{\R^d}x_1^2\,\eu^{-\norm x^2}\dd x =
\frac{\sqrt\pi}2\;\pi^{(d-1)/2} = \frac{\pi^{d/2}}2,
\qquad
\int_{\R^d}\norm x^2\eu^{-\norm x^2}\dd x =
d\cdot\frac{\pi^{d/2}}2
\]
(सममिति से $\norm x^2 = \sum_ix_i^2$ $d$ बराबर पदों का योगदान देता है --- यही संगति
जाँच है)। मानकीकृत \omterm{def:b3:measure:measure}{माप} $\pi^{-d/2}\eu^{-\norm x^2}\dd x$ के लिए द्वितीय आघूर्ण $\frac d2$ है।
\end{solution}

\begin{solution}{exo:b3:product:11}
(क) $\{y > 0\}$ से प्रतिच्छेदित $\Phi(x, y) = f(x) -
y$ के लिए $H = \Phi^{-1}(\intoo0\infty)$: \omterm{def:b3:lebesgue:measurable}{मापनीय}, क्योंकि $(x, y)
\mapsto f(x)$
और $(x,y)\mapsto y$ गुणनफल पर \omterm{def:b3:lebesgue:measurable}{मापनीय} हैं (प्रक्षेपों के साथ संयोजन)। $H$ का
$x$-अनुभाग $\intoo0{f(x)}$ है, जिसका \omterm{def:b3:measure:measure}{माप} $f(x)$ है: टोनेली अनुभागों का समाकलन
करती है,
\[
\lambda_{d+1}(H) = \int_{\R^d}\lambda_1\bigl(\intoo0{f(x)}
\bigr)\,\dd x = \int_{\R^d}f\,\dd\lambda_d .
\]

(ख) आलेख $\{(x,y) : y \geq f(x)\} \cap \{y \leq
f(x)\}$ है, जो \omterm{def:b3:lebesgue:measurable}{मापनीय} है; उसके $x$-अनुभाग एकल बिंदु हैं,
जिनका \omterm{def:b3:measure:measure}{माप} $0$ है: अतः टोनेली $\lambda_{d+1}(\text{आलेख}) =
\int 0 = 0$ देती है।

(ग) $S^{d-1}$ $\R^{d-1}$ के इकाई गोले पर दोनों आलेखों $y =
\pm\sqrt{1 - \abs{x'}^2}$ का संघ है (अंतिम
निर्देशांक अलग करके): अर्थात् (ख) से दो अकिंचन समुच्चयों का संघ,
इसलिए अकिंचन।
\end{solution}

\begin{solution}{exo:b3:product:12}
$\intoo01^2$ पर ऋणेतर पदों के साथ $\frac1{1 - xy} = \sum_{n\geq0}(xy)^n$: टोनेली पदशः समाकलन की अनुमति देती
है,
\[
\iint\frac{\dd x\,\dd y}{1 - xy}
= \sum_{n\geq0}\Bigl(\int_0^1x^n\dd x\Bigr)
\Bigl(\int_0^1y^n\dd y\Bigr)
= \sum_{n\geq0}\frac1{(n+1)^2} = \zeta(2) .
\]
एकांतर स्थिति के लिए $\frac1{1 + xy} =
\sum_n(-1)^n(xy)^n$ कोई धनात्मक श्रेणी नहीं है; पर \emph{निरपेक्ष}
श्रेणी का समाकल $\zeta(2) <
\infty$ है, अतः फुबिनी (अब \omterm{def:b3:lebesgue:l1}{समाकलनीयता} स्थापित) लागू
होती है: $\iint\frac{\dd x\dd y}{1 + xy} =
\sum_n\frac{(-1)^n}{(n+1)^2}$। श्रेणी सर्वसमिका:
\[
\sum_{n\geq1}\frac{(-1)^{n-1}}{n^2}
= \sum_{n\geq1}\frac1{n^2} - 2\sum_{k\geq1}\frac1{(2k)^2}
= \zeta(2) - \frac{\zeta(2)}2 = \frac{\zeta(2)}2 .
\]
$\zeta(2) = \frac{\pi^2}6$ (\cref{pb:b3:spectral:1}) के साथ: समाकल $\frac{\pi^2}6$ और $\frac{\pi^2}{12}$ हैं। पहले परिकलन में
टोनेली की धनात्मकता ही सब कुछ थी --- अदला-बदली से पहले किसी \omterm{def:b3:lebesgue:l1}{समाकलनीयता}
जाँच की आवश्यकता नहीं; और दूसरे में निरपेक्ष श्रेणी की धनात्मकता ही
\omterm{def:b3:lebesgue:l1}{समाकलनीयता} \emph{प्रमाणित} करती है ताकि फुबिनी चिह्नयुक्त श्रेणी पर
चल सके।
\end{solution}

\begin{solution}{pb:b3:product:1}
\textbf{1.} $x = t + \sqrt t\,u$ के साथ ($\dd x = \sqrt
t\,\dd u$; जैसे-जैसे $u$ $\intoo{-\sqrt t}\infty$ पर चलता है
$x$ $\intoo0\infty$ पर चलता है):
\[
\Gamma(t{+}1) = \int_0^\infty x^t\eu^{-x}\dd x
= t^t\eu^{-t}\sqrt t\int_{-\sqrt t}^{\infty}
\exp\Bigl(t\ln\Bigl(1 + \frac u{\sqrt t}\Bigr) - \sqrt
t\,u\Bigr)\dd u,
\]
क्योंकि $x^t = t^t\exp\bigl(t\ln(1 + u/\sqrt t)\bigr)$ और $\eu^{-x} = \eu^{-t}\eu^{-\sqrt tu}$।

\textbf{2.} स्थिर $u$ और $t \to \infty$ के लिए: $t\ln(1 +
u/\sqrt t) - \sqrt tu = t\bigl(\frac u{\sqrt t} -
\frac{u^2}{2t} + o(\frac1t)\bigr) - \sqrt tu = -\frac{u^2}2 +
o(1)$: $g_t(u) \to \eu^{-u^2/2}$।

\textbf{3.} $\intoc{-1}1$ पर $\psi_1(h) = \varphi(h) + \frac{h^2}4$ रखिए: $\psi_1(0) = 0$ और $\psi_1'(h) = \frac1{1+h} -
1 + \frac h2 = \frac{h(h-1)}{2(1+h)}$, जो $\intoc{-1}0$ पर $\geq 0$
और $\intcc01$ पर $\leq 0$ है: $\psi_1 \leq 0$, अर्थात् वहाँ $\varphi(h) \leq -h^2/4$। $\intco1\infty$ पर $\psi_2(h) =
\varphi(h) + ch$,
$c = 1 - \ln2$ रखिए: $\psi_2(1)
= \ln2 - 1 + c = 0$ और $\psi_2'(h) = c - \frac h{1+h} \leq c -
\frac12 < 0$: $h \geq 1$ के लिए $\varphi(h) \leq -ch$। अब $t
\geq 1$ के लिए:
यदि $\abs u \leq \sqrt t$ हो, तो $g_t(u) =
\eu^{t\varphi(u/\sqrt t)} \leq \eu^{-t(u/\sqrt t)^2/4} =
\eu^{-u^2/4}$; और यदि $u \geq \sqrt t$ हो, तो $t\,\varphi(u/\sqrt t) \leq
-ct\cdot\frac u{\sqrt t} = -c\sqrt t\,u \leq -cu$ (क्योंकि $t \geq
1$),
अतः $g_t(u) \leq \eu^{-cu}$। इसलिए $g_t \leq \eu^{-u^2/4} +
\eu^{-cu}\mathbf 1_{u>0}$, जो \omterm{def:b3:lebesgue:l1}{समाकलनीय} और $t \geq
1$ से स्वतंत्र है।

\textbf{4.} प्रभावी अभिसरण प्रमेय: $\int_\R g_t(u)\dd u \to
\int_\R\eu^{-u^2/2}\dd u = \sqrt{2\pi}$ (\cref{ex:b3:product:polar} और साथ में मापन
$u\mapsto\sqrt2\,u$)। प्रश्न 1 के साथ:
\[
\Gamma(t + 1) \sim \sqrt{2\pi t}\;\Bigl(\frac
t\eu\Bigr)^t,\qquad
n! \sim \sqrt{2\pi n}\,\Bigl(\frac n\eu\Bigr)^n .
\]

\textbf{5.} \cref{exo:b3:product:8} से $W_{2n} =
\frac\pi2\,\frac{(2n)!}{4^n(n!)^2} =
\frac\pi2\,4^{-n}\binom{2n}n$। स्टर्लिंग:
\[
\binom{2n}{n} = \frac{(2n)!}{(n!)^2}
\sim \frac{\sqrt{4\pi n}\,(2n/\eu)^{2n}}
{2\pi n\,(n/\eu)^{2n}} = \frac{4^n}{\sqrt{\pi n}} .
\]
तब $W_{2n} \sim \frac12\sqrt{\pi/n}$, जो पुनरावृत्ति $W_n = \frac{n-1}nW_{n-2}$ के संगत है (जो $W_n \sim
W_{n-2}$ अनिवार्य कर देती
है, और $W_nW_{n-1}\cdot n = \frac\pi2$ --- शास्त्रीय वालिस संबंध --- के साथ $W_n \sim
\sqrt{\pi/(2n)}$ दे देती है;
दोनों अनंतस्पर्शी सहमत हैं)।

\textbf{6.} $\Gamma(\frac d2 + 1) \sim \sqrt{2\pi\frac
d2}\,(\frac d{2\eu})^{d/2}$, अतः
\[
v_d = \frac{\pi^{d/2}}{\Gamma(\frac d2 + 1)}
\sim \frac{1}{\sqrt{\pi d}}\Bigl(\frac{2\pi\eu}
d\Bigr)^{d/2} \longrightarrow 0
\]
अति-ज्यामितीय रूप से ($d > 2\pi\eu \approx 17$ के लिए प्रत्येक गुणक $< 1$ और सिकुड़ता
हुआ)। संख्यात्मक रूप से $v_1 = 2$, $v_2 \approx
3.14$, $v_3 \approx 4.19$, $v_4 \approx 4.93$, $v_5 \approx
5.26$, $v_6 \approx 5.17$:
उच्चिष्ठ $d = 5$ पर है।

\textbf{7.} $k = \frac n2 + \frac{s\sqrt n}2$ के साथ (पूर्णांक, $n$ सम, $s$ स्थिर): $2^{-n}\binom nk =
2^{-n}\frac{n!}{k!(n-k)!}$
में लघुगणक लीजिए और तीनों क्रमगुणितों पर स्टर्लिंग लगाइए। $\varepsilon = s/\sqrt n$ के
साथ $k = \frac n2(1 + \varepsilon)$, $n - k =
\frac n2(1 - \varepsilon)$ लिखने पर:
\[
\ln\Bigl(2^{-n}\binom nk\Bigr)
= -\frac n2\bigl[(1{+}\varepsilon)\ln(1{+}\varepsilon) +
(1{-}\varepsilon)\ln(1{-}\varepsilon)\bigr]
+ \frac12\ln\frac{2}{\pi n(1 - \varepsilon^2)} + o(1),
\]
और कोष्ठक $\varepsilon^2 + O(\varepsilon^4) =
\frac{s^2}n + O(n^{-2})$ है: अतः प्रदर्शन $o(1)$ तक $-\frac{s^2}2 +
\frac12\ln\frac2{\pi n}$ की ओर जाता है, अर्थात्
\[
2^{-n}\binom nk \sim \sqrt{\frac{2}{\pi
n}}\;\eu^{-s^2/2} :
\]
अर्थात् सिक्का-उछाल का गाउसीय परिच्छेद, मात्रात्मक रूप में ---
दे म्वाव्र--लाप्लास का स्थानीय रूप, जिसे \cref{ch:b3:clt} में वैश्विक बनाया
जाएगा।

\textbf{8.} (क) प्रभावी अभिसरण प्रमेय प्रश्न 3 के प्रभावी फलन का
उपयोग करते हुए प्रश्न 2 की बिंदुशः सीमा को समाकलों के अभिसरण में
बदल देती है। (ख) गाउसीय समाकल सीमा $\int\eu^{-u^2/2} = \sqrt{2\pi}$ का मूल्यांकन कर देता है
--- स्टर्लिंग का अचर $\sqrt{2\pi}$ \emph{ही} गाउसीय समाकल है। (ग) प्रतिस्थापन
$x = t + \sqrt tu$ कोई आफ़ीन चर परिवर्तन है: स्थानांतरण निश्चरता और \omterm{def:b3:measure:lebesgueouter}{लेबेग माप} का
मापन नियम (विमा $1$ में \cref{thm:b3:product:linearchange})।

\textbf{9.} दोनों पदों का प्रसार कीजिए:
\[
d_n - d_{n+1} = \ln\frac{n!}{(n+1)!} + \Bigl(n +
\frac32\Bigr)\ln(n+1) - \Bigl(n + \frac12\Bigr)\ln n - 1
= \Bigl(n + \frac12\Bigr)\ln\frac{n+1}n - 1,
\]
जिसमें पद $-\ln(n+1)$ और $(n + \frac32)\ln(n+1)$ मिलकर $(n + \frac12)\ln(n+1)$ बन जाते हैं।

\textbf{10.} $t = \frac1{2n+1}$ के लिए: $\frac{1+t}{1-t} =
\frac{2n+2}{2n} = \frac{n+1}n$, और $n + \frac12 =
\frac1{2t}$; विषम श्रेणी $\ln\frac{1+t}{1-t} =
2\sum_{k\geq0}\frac{t^{2k+1}}{2k+1}$
\[
\Bigl(n + \frac12\Bigr)\ln\frac{n+1}n
= \sum_{k\geq0}\frac{t^{2k}}{2k+1}
= 1 + \frac{t^2}3 + \frac{t^4}5 + \cdots
\]
देती है। $1$ घटाइए। निचला परिबंध: अकेला पहला पद, $\frac{t^2}3 = \frac1{3(2n+1)^2}$। ऊपरी
परिबंध: सभी हरों को $3$ तक नीचे लाइए और ज्यामितीय श्रेणी का योग
लीजिए: $\frac{t^2}{3(1 - t^2)} = \frac1{3((2n+1)^2 - 1)} =
\frac1{12n(n+1)} = \frac1{12n} - \frac1{12(n+1)}$।

\textbf{11.} ऊपरी परिबंध का $n$ से $\infty$ तक योग ($d_m \to 0$ के साथ):
$d_n < \frac1{12n}$। निचले परिबंध के लिए: $\frac1{12m+1} - \frac1{12(m+1)+1} =
\frac{12}{(12m+1)(12m+13)}$, और
\begin{align*}
\frac1{3(2m+1)^2} > \frac{12}{(12m+1)(12m+13)}
&\iff (12m+1)(12m+13) > 36(2m+1)^2 \\
&\iff 168m + 13 > 144m + 36,
\end{align*}
जो $m \geq 1$ के लिए सत्य है। इस दूरबीनी लघुकारक का योग: $d_n > \frac1{12n+1}$। चरघातांक
लेने पर $n!$ का शास्त्रीय परिबंधन मिल जाता है।

\textbf{12.} (क) बड़े $n$ के लिए सापेक्ष त्रुटि $= \eu^{d_n} - 1 <
\eu^{1/(12n)} - 1 < \frac{1.1}{12n}$; $12n \geq 1.1\cdot10^6$ होते
ही $<
10^{-6}$, और कहा गया $n \geq 83\,334$ पर्याप्त है ($\frac1{12n} \leq 10^{-6}$ पहले से ही इसे निहित
करता है)। (ख) $\log_{10}(100!) =
\frac12\log_{10}(200\pi) + 200 - 100\log_{10}\eu +
d_{100}\log_{10}\eu = 1.39906 + 200 - 43.42945 + 0.00036
\approx 157.96997$, अतः $100! \approx 10^{0.96997} \cdot
10^{157} = 9.333\cdot10^{157}$; और प्रत्याभूत खिड़की $(\eu^{1/1201}, \eu^{1/1200})$ की चौड़ाई
सापेक्ष रूप में $10^{-6}$ से कम है --- अर्थात् एक ``अनंतस्पर्शी'' सूत्र
जो $n =
100$ पर परिशुद्धता का यंत्र है।

\textbf{13.} $\sin^n = \sin^{n-2} - \sin^{n-2}\cos^2$ लिखिए और दूसरे पद का खंडशः समाकलन कीजिए ($u = \cos\theta$,
$\dd v = \sin^{n-2}\cos\theta\,\dd\theta$, $v =
\frac{\sin^{n-1}}{n-1}$):
\[
\int_0^{\pi/2}\sin^{n-2}\cos^2 =
\Bigl[\cos\theta\,\frac{\sin^{n-1}\theta}{n-1}\Bigr]_0^{\pi/2}
+ \frac1{n-1}\int_0^{\pi/2}\sin^n = \frac{W_n}{n-1} .
\]
अतः $W_n = W_{n-2} - \frac{W_n}{n-1}$, अर्थात् $W_n =
\frac{n-1}nW_{n-2}$। $W_0 = \frac\pi2$, $W_1 = 1$ से:
\[
W_{2n} = \frac{(2n-1)!!}{(2n)!!}\cdot\frac\pi2
= \frac\pi2\binom{2n}n4^{-n},
\qquad
W_{2n+1} = \frac{(2n)!!}{(2n+1)!!}
= \frac{4^n}{(2n+1)\binom{2n}n},
\]
जिसमें द्विक क्रमगुणितों को $(2n)!! = 2^nn!$ और $(2n-1)!! = \frac{(2n)!}{2^nn!}$ से बदला गया है। अंत में
पुनरावृत्ति से $nW_nW_{n-1} =
(n-1)W_{n-1}W_{n-2}$: अचर, जो $1\cdot W_1W_0 = \frac\pi2$ के बराबर है।

\textbf{14.} $W_{2n+1} \leq W_{2n} \leq W_{2n-1}$ ($\sin^n$ की बिंदुशः एकदिष्टता) और $\frac{W_{2n-1}}{W_{2n+1}} =
\frac{2n+1}{2n} \to 1$ $\frac{W_{2n}}{W_{2n+1}} \to
1$ को
दबा देते हैं। $W_{2n}W_{2n+1} = \frac{\pi}{2(2n+1)}$ (प्रश्न 13) के साथ मिलाकर: $W_{2n}^2 \sim \frac\pi{4n}$, अतः $W_{2n} \sim
\frac12\sqrt{\frac\pi n}$ और
$\binom{2n}n4^{-n} =
\frac2\pi W_{2n} \sim \frac1{\sqrt{\pi n}}$।

\textbf{15.} प्रश्न 9--10 ने अचर के मान का उपयोग कभी नहीं किया:
$e_n = \ln n! - (n+\frac12)\ln n + n$ के साथ अंतर $e_n - e_{n+1}$ $\bigl(0, \frac1{12n} -
\frac1{12(n+1)}\bigr)$ में होते हैं, अतः $(e_n)$ घटता है जबकि
$(e_n -
\frac1{12n})$ बढ़ता है: निकटवर्ती अनुक्रम, जो किसी उभयनिष्ठ $\ell$ पर अभिसरित
होते हैं। अतः $n! \sim K n^{n+1/2}\eu^{-n}$, $K =
\eu^\ell$।

\textbf{16.} अज्ञात-अचर वाली स्टर्लिंग को केंद्रीय द्विपद में
प्रतिस्थापित करने पर:
\[
\binom{2n}n4^{-n}\sqrt n \sim
\frac{K\,(2n)^{2n+1/2}\eu^{-2n}}{\bigl(K\,n^{n+1/2}
\eu^{-n}\bigr)^2}\,4^{-n}\sqrt n
= \frac{2^{2n}\sqrt{2}\,K\,n^{2n+1/2}}{K^2\,n^{2n+1}}
\,4^{-n}\sqrt n = \frac{\sqrt2}{K},
\]
और प्रश्न 14 $\frac{\sqrt2}K = \frac1{\sqrt\pi}$ अनिवार्य कर देता है: $K = \sqrt{2\pi}$। इस प्रकार भाग III
और IV स्टर्लिंग को शून्य से पुनः सिद्ध कर देते हैं; और उसे भाग I
की सर्वसमिका में डालने पर $\int_\R\eu^{-u^2/2}\dd u = \sqrt{2\pi}$ का मूल्यांकन बिना ध्रुवीय निर्देशांकों
के हो जाता है: वालिस और गाउस एक-दूसरे को सहारा देते हैं।

\textbf{17.} $j \geq 0$ के लिए $\frac{\binom{2n}{n+j}}{\binom{2n}n} =
\frac{(n!)^2}{(n+j)!\,(n-j)!} =
\prod_{i=1}^{j}\frac{n-i+1}{n+i}$ (और सममिति से $j < 0$ के लिए भी)।
$1 \leq i
\leq j \leq K\sqrt n$ के साथ लघुगणक लेने पर:
\[
\ln\frac{n-i+1}{n+i} = \ln\Bigl(1 - \frac{i-1}n\Bigr) -
\ln\Bigl(1 + \frac in\Bigr) = -\frac{2i-1}{n} +
O\Bigl(\frac{i^2}{n^2}\Bigr),
\]
और $\sum_{i\leq j}(2i - 1) = j^2$, जबकि त्रुटि का योग $O(j^3/n^2) = O(n^{-1/2})$ है: अतः एकसमान रूप से $\exp\bigl(-\frac{j^2}n + O(n^{-1/2})\bigr)$।

\textbf{18.} $\eu^{-n}\frac{n^n}{n!} \sim \eu^{-n}
\frac{n^n}{\sqrt{2\pi n}\,n^n\eu^{-n}} =
\frac1{\sqrt{2\pi n}}$। माध्य $n$ वाले किसी प्वासों चर का मानक विचलन
$\sqrt n$ होता है, और $\frac1{\sqrt{2\pi n}}$ ठीक गाउसीय शिखर ऊँचाई $\frac1{\sigma\sqrt{2\pi}}$ है: अर्थात्
स्थानीय केंद्रीय सीमा प्रमेय, बहुलक पर पहले ही झलक जाती हुई।

\textbf{19.} $a \in \intoo01$ के लिए \cref{pb:b3:lebesgue:1} की ढाल प्रमेयिका (वहाँ प्रश्न 14),
$n$ के चारों ओर उत्तल $\log\Gamma$ पर लगाई गई, $(n-1)^a \leq
\frac{\Gamma(n+a)}{\Gamma(n)} \leq n^a$ देती है: अतः $n^a$
से अनुपात $(1 - \frac1n)^a \to 1$ से दब जाता है। $a = m + a'$ ($m \in \N$, $a' \in \intco01$) के लिए: $\Gamma(n+a) = (n + a - 1)
\cdots(n + a')\Gamma(n + a')$,
और $m$ गुणकों में से प्रत्येक $n(1 + O(\frac1n))$ है: आकलनों को गुणा कीजिए।

\textbf{20.} पुनरावृत्ति $v_d = \frac{2\pi}dv_{d-2}$ ($\Gamma(\frac d2 + 1) = \frac d2\Gamma(\frac d2)$ से) $v_1 = 2$, $v_2 = \pi$ से
\[
v_3 = \frac{4\pi}3,\quad v_4 = \frac{\pi^2}2,\quad
v_5 = \frac{8\pi^2}{15},\quad v_6 = \frac{\pi^3}6,\quad
v_7 = \frac{16\pi^3}{105}.
\]
देती है। अनुपात $\frac{2\pi}d$ ठीक $d \leq 6$ के लिए $1$ से अधिक है, अतः प्रत्येक
सम-विषमता पहले बढ़ती फिर घटती है; संख्यात्मक रूप से $v_4
\approx 4.93$, $v_5 \approx 5.26$,
$v_6 \approx 5.17$: अतः समग्र उच्चिष्ठ $d = 5$ है। जनक फलन: $v_{2k} =
\frac{\pi^k}{k!}$, अतः $\sum_kv_{2k}x^{2k} =
\eu^{\pi x^2}$ ---
अर्थात् सम-विमीय सभी गोला आयतन एक ही चरघातांक में लिपटे हुए, और
$v_d$ के लिए तत्काल अति-ज्यामितीय क्षय आकलन।

\textbf{21.} अंश और हर में स्टर्लिंग, $k =
\alpha n$ के साथ:
\[
\binom n{\alpha n} \sim
\frac{\sqrt{2\pi n}\,n^n}
{\sqrt{2\pi\alpha n}\,(\alpha n)^{\alpha n}\,
\sqrt{2\pi(1-\alpha)n}\,((1-\alpha)n)^{(1-\alpha)n}}
= \frac{\eu^{nH(\alpha)}}{\sqrt{2\pi\alpha(1-\alpha)n}},
\]
क्योंकि $n^n/(\alpha n)^{\alpha n}((1-\alpha)n)^{(1-\alpha)n} =
\alpha^{-\alpha n}(1-\alpha)^{-(1-\alpha)n} =
\eu^{nH(\alpha)}$ ($n$ की घातें निरस्त हो जाती हैं: $\alpha n +
(1-\alpha)n = n$), और $\eu^{-n}$
भी वैसे ही निरस्त हो जाते हैं। $\alpha = \frac12$ पर: $H = \ln2$ और पूर्व-गुणक $\sqrt{2/(\pi n)}$
है --- फिर से प्रश्न 5। $H$ की कड़ी अवतलता (उसका दूसरा अवकलज $-\frac1{\alpha(1-\alpha)} < 0$)
उसका उच्चिष्ठ $\ln 2$ केवल $\alpha = \frac12$ पर रखती है: $\alpha \neq \frac12$ के लिए $\binom n{\alpha n}2^{-n} \approx
\eu^{-n(\ln2 - H(\alpha))}$
चरघातांकी रूप से क्षय होता है --- सिक्का-उछालों के बारे में प्रत्येक
संकेंद्रण कथन के पीछे यही संयोजनात्मक इंजन है।

\textbf{22.} $s_{d-1} = dv_d$ और प्रश्न 20 से:
\[
s_0 = 2,\ \ s_1 = 2\pi,\ \ s_2 = 4\pi,\ \ s_3 = 2\pi^2,\ \
s_4 = \frac{8\pi^2}3,\ \ s_5 = \pi^3,\ \ s_6 =
\frac{16\pi^3}{15},
\]
संख्यात्मक रूप से $2,\ 6.28,\ 12.57,\ 19.74,\ 26.32,\ 31.01,\
33.07$; और $s_7 = \frac{\pi^4}3 \approx 32.47 < s_6$: अतः उच्चिष्ठ $6$-गोलक है।
पुनरावृत्ति $s_{d+1} =
(d+2)\,v_{d+2} = (d+2)\,\frac{2\pi}{d+2}\,v_d = 2\pi v_d =
\frac{2\pi}d\,s_{d-1}$ आयतनों जैसी ही $\frac{2\pi}d$-चालित वृद्धि और अति-ज्यामितीय
गिरावट दिखाती है: सातवीं विमा के बाद गोलक किसी भी ज्यामितीय अनुक्रम
से तेज़ी से सिकुड़ जाते हैं।

\textbf{23.} प्रश्न 11 ठीक $\frac1{12n+1} < d_n <
\frac1{12n}$ कहता है, और
\[
\frac1{12n} - \frac1{12n+1} = \frac1{12n(12n+1)} =
O\Bigl(\frac1{n^2}\Bigr),
\]
अतः $d_n = \frac1{12n} + O(\frac1{n^2})$ और $\eu^{d_n} = 1 +
\frac1{12n} + O(\frac1{n^2})$; $\sqrt{2\pi n}(n/\eu)^n$ से गुणा करने पर संशोधित सूत्र मिल जाता
है। $n =
10$ पर: $\sqrt{20\pi}\,(10/\eu)^{10} = 7.92665 \times 453999.3
\approx 3\,598\,696$, जो $30\,104$ कम है (सापेक्ष त्रुटि $8.3\cdot10^{-3}$); और $1 + \frac1{120}$
से गुणा करने पर $3\,628\,685$ मिलता है, जो $115$ कम है (सापेक्ष त्रुटि
$3.2\cdot10^{-5}$)। स्वयं यह परिबंधन $10!$ को $3\,598\,696\,\eu^{1/121} \approx 3\,628\,559$ और $3\,598\,696\,\eu^{1/120} \approx 3\,628\,808$ के बीच बाँध देता है
--- ऊपरी परिबंध सात अंकों में आठ इकाई से चूकता है।

\textbf{24.} भाग I का प्रतिस्थापन $x = t + \sqrt t\,u$, कर्तित समाकल पर लगाया गया,
\[
\int_0^{t} x^{t}\eu^{-x}\,\dd x
= t^{t}\eu^{-t}\sqrt t\int_{-\sqrt t}^{0} g_t(u)\,\dd u ,
\]
देता है, जिसमें परास $0 \leq x \leq t$ $-\sqrt t \leq u \leq
0$ बन जाता है। प्रश्न 3 का प्रभावी
फलन $g_t\mathbf 1_{u < 0}$ को भी ढक लेता है, अतः प्रभावी अभिसरण
\[
\int_{-\sqrt t}^{0}g_t(u)\,\dd u \longrightarrow
\int_{-\infty}^{0}\eu^{-u^2/2}\dd u = \frac{\sqrt{2\pi}}2 ,
\]
दे देता है, जबकि प्रश्न 4 $\Gamma(t+1) \sim
t^t\eu^{-t}\sqrt t\,\sqrt{2\pi}$ देता है। अनुपात $\frac12$ की ओर जाता
है। प्रायिकता की दृष्टि से: बड़े आकार-प्राचल वाला कोई गामा यादृच्छिक
चर अपने बहुलक के दोनों ओर अनंतस्पर्शी रूप से आधा-आधा द्रव्यमान रखता
है --- अर्थात् केंद्रीय सीमा प्रमेय की सममिति, एक ही प्रतिस्थापन
से पढ़ ली गई।

\textbf{25.} मान लीजिए $\lambda = \frac{\alpha}{1-\alpha} \in
\intoc01$। $k \leq \lfloor\alpha n\rfloor \leq \alpha n$ के लिए हमारे पास $\lambda^{k - \alpha n} \geq 1$ है, अतः
\[
\sum_{k=0}^{\lfloor\alpha n\rfloor}\binom nk
\leq \lambda^{-\alpha n}\sum_{k=0}^{n}\binom nk\lambda^{k}
= \Bigl(\lambda^{-\alpha}(1 + \lambda)\Bigr)^{n},
\]
और $\lambda = \frac\alpha{1-\alpha}$ के साथ:
\[
\lambda^{-\alpha}(1+\lambda)
= \alpha^{-\alpha}(1-\alpha)^{\alpha}\cdot\frac1{1-\alpha}
= \alpha^{-\alpha}(1-\alpha)^{-(1-\alpha)}
= \eu^{H(\alpha)} .
\]
अनुकूलतमता: $\lambda > 0$ पर $f(\lambda) = -\alpha\ln\lambda +
\ln(1+\lambda)$ का न्यूनीकरण करने पर समीकरण $f'(\lambda)
= -\frac\alpha\lambda + \frac1{1+\lambda} = 0$ का
अद्वितीय हल $\lambda = \frac\alpha{1-\alpha}$ है, जो निम्निष्ठ है क्योंकि $f'' > 0$ --- अर्थात्
चरघातांकी-नमन (चेर्नोफ) वाला चयन। मेल: प्रश्न 21 से अकेला पद $k =
\lfloor\alpha n\rfloor$
पहले से ही $\eu^{nH(\alpha)}/\sqrt{2\pi\alpha(1-\alpha)n}$ कोटि का है, अतः
\[
\frac{\eu^{nH(\alpha)}}{C\sqrt n} \leq
\sum_{k\leq\alpha n}\binom nk \leq \eu^{nH(\alpha)} :
\]
अर्थात् दर $H(\alpha)$ यथार्थ है, और समूचा योग अपने सबसे बड़े पद पर अधिक
से अधिक $\sqrt n$ गुणक की क़ीमत लेता है। $2^n$ से भाग देने पर यह निष्पक्ष
सिक्के का पुच्छ परिबंध $\P(S_n \leq \alpha n) \leq
\eu^{-n(\ln2 - H(\alpha))}$ है --- अर्थात् एक पंक्ति में \omterm{def:b3:measure:measure}{माप} का
संकेंद्रण।
\end{solution}
